% docs-maximize-roi — HBR-magazine house-style template (LaTeX, pdflatex).
% Matches ~/dev2/aicostoptimize/docs/ROI/hbr-cfo-article.tex and
% ~/dev2/aicostmemory/docs/ROI/mindspace-hbr.tex. Two-column magazine layout:
% red kicker, Huge sans title, gray dek, byline, "Idea in Brief" box, red
% pull-quote, full-width EXHIBIT (image OR native chart), footer.
% Replace every % FILL: block. Compiles with: pdflatex <slug>-hbr.tex
\documentclass[10pt]{article}

\usepackage[margin=0.95in,top=0.85in]{geometry}
\usepackage{amsmath}
\usepackage{booktabs}
\usepackage{xcolor}
\usepackage{graphicx}   % for an image exhibit (e.g. a Mermaid quadrantChart -> PNG)
\usepackage{tikz}       % optional: for a native chart exhibit instead of an image
\usepackage{pgfplots}
\pgfplotsset{compat=1.17}
\usepackage[hidelinks]{hyperref}
\usepackage{microtype}
\usepackage{helvet}     % sans for headings/labels
\usepackage{sectsty}
\usepackage{multicol}
\usepackage{mdframed}
\usepackage{enumitem}

\definecolor{hbrred}{RGB}{173,31,45}
\definecolor{hbrdark}{RGB}{35,31,32}
\definecolor{hbrgray}{RGB}{100,100,105}
\definecolor{hbrbg}{RGB}{245,242,236}

\allsectionsfont{\sffamily\color{hbrdark}}
\newcommand{\dek}[1]{{\large\sffamily\color{hbrgray} #1\par}}
\newenvironment{ideabrief}
  {\begin{mdframed}[backgroundcolor=hbrbg,linecolor=hbrred,linewidth=2pt,
    topline=false,bottomline=false,rightline=false,
    innerleftmargin=10pt,innerrightmargin=10pt,innertopmargin=8pt,innerbottommargin=8pt]
   \sffamily\small}
  {\end{mdframed}}
\newcommand{\pullquote}[1]{%
  \par\vspace{6pt}
  {\noindent\color{hbrred}\rule{\linewidth}{0.6pt}}\par\vspace{4pt}
  {\noindent\large\sffamily\itshape\color{hbrred} ``#1''}\par\vspace{4pt}
  {\noindent\color{hbrred}\rule{\linewidth}{0.6pt}}\par\vspace{6pt}}

\setlength{\parskip}{4pt}
\setlength{\parindent}{0pt}
\setlength{\columnsep}{22pt}

\begin{document}

% --- kicker: section label (left) + status (right) -------------------------
{\sffamily\footnotesize\color{hbrred} FILL:SECTION \,\textbullet\, FILL:TOPIC \hfill FILL: INTERNAL DRAFT v0.1 --- STEALTH, NOT FOR DISTRIBUTION}

\vspace{4mm}
% --- title: keep 2-3 short lines; use \\ for the line breaks ---------------
{\Huge\sffamily\bfseries\color{hbrdark} FILL: Line one\\ FILL: line two\\ FILL: line three.\par}

\vspace{3mm}
\dek{FILL: the deck — one or two sentences that reframe the reader's problem, in gray sans.}

\vspace{2mm}
{\sffamily\small by \textbf{FILL: Author}, FILL: Company \hfill FILL: Month Year}

\vspace{4mm}
\begin{ideabrief}
\textbf{THE IDEA IN BRIEF}

\textbf{The problem.} % FILL

\textbf{Why the usual fixes fail.} % FILL

\textbf{The move.} % FILL (the mechanism from the ROI analysis, in business terms)

\textbf{The payoff.} % FILL (the compounding / the moat; end on the reframe line)
\end{ideabrief}

\vspace{3mm}
\begin{multicols}{2}

\section*{FILL: opening section — the pain the reader feels}
% FILL: 2 short paragraphs in the executive's language. End on the sharp observation.

\section*{FILL: why the obvious fixes fall short}
% FILL: dispatch the alternatives; each optimizes the easy thing and ignores the valuable one.

\section*{FILL: the reframe — <mechanism> in business terms}
% FILL: the load-bearing properties, each with "remove it and you become <inferior alternative>."

\pullquote{FILL: the single most quotable line — the thesis in one sentence.}

% FILL: one line teeing up the Exhibit.

\end{multicols}

% --- full-width EXHIBIT: an image (Mermaid -> PNG) OR a native chart --------
\begin{figure}[h]
\centering
% OPTION A — image exhibit (recommended for a 2x2 quadrant; render via Mermaid).
% Uncomment and set the filename once the PNG exists:
%   \includegraphics[width=0.62\textwidth]{FILL-exhibit.png}
% OPTION B — native chart exhibit: a tikzpicture/pgfplots axis (see aicostoptimize
% hbr-cfo-article.tex for a savings-curve example).
% The framebox below is a placeholder so the template compiles as-is — replace it.
\framebox[0.62\textwidth]{\rule{0pt}{4cm}\sffamily\color{hbrgray} FILL: exhibit (image or chart)}

\vspace{1mm}
\begin{minipage}{0.86\textwidth}
\footnotesize\sffamily\color{hbrgray}
\textbf{EXHIBIT} \,% FILL: caption. If projections, label them projections and state assumptions.
\end{minipage}
\end{figure}

\begin{multicols}{2}

\section*{FILL: the flywheel / why it compounds}
% FILL: why the advantage grows with use.

\section*{FILL: what the numbers actually say — in business terms}
% FILL: the honest ranking, most->least defensible; de-rank the oversold headline out loud.
\begin{enumerate}[leftmargin=1.4em, itemsep=2pt]
\item \textbf{FILL: the provable near-term dollar.} % ...
\item \textbf{FILL: the durable moat.} % ...
\item \textbf{FILL: the compounding flywheel.} % ...
\item \textbf{FILL: the adoption lever.} % ...
\item \textbf{FILL: the oversold headline (de-ranked).} % real but weakest; don't lead with it.
\end{enumerate}

\section*{What to do Monday}
\begin{itemize}[leftmargin=1.2em, itemsep=2pt]
\item \textbf{FILL.} % concrete action
\item \textbf{FILL.} % concrete action
\item \textbf{FILL.} % concrete action
\item \textbf{FILL: turn measurement into action.} % the compounding, not just reporting
\end{itemize}

% FILL: closing paragraph — the memorable last line.

\end{multicols}

\vspace{2mm}
\noindent\rule{\textwidth}{0.4pt}

\footnotesize\sffamily\color{hbrgray}
\textbf{FILL: Author} % FILL: one-line bio; reference the companion academic paper; note that
% the strongest claim is architectural, not numerical; cite real sources for any market stats.

\end{document}
